\section{Morley's Original Categoricity-in-Power Argument}
\label{sec:p03-n-a-006-01}

Morley's 1965 argument proves more than transfer of categoricity: if a
countable theory is categorical in one uncountable cardinal, every
uncountable model is saturated. Its machinery consists of a derivative
of all parameter type spaces together, a Ramsey construction of models
with few types realized over countable sets, and prime models arranged
compatibly along chains. These are the original five sections of
\emph{Categoricity in Power}, including its three rank examples and
final rank bound \cite[pp.~514--538]{Morley1965}. Strongly minimal
pregeometry and prime-over-basis coordinatization are later developments;
they are not substituted for this proof.

\subsection{Normalization and the category of parameter systems}
\label{sec:p03-morley-normalization}

\begin{hbconvention}{def:p03-morley-standing}{The original standing assumptions}
Let \(\Sigma\) be a satisfiable complete theory in an at most countable
first-order language, with an infinite model. First replace any
functions and constants by their total single-valued
graph relations and singleton relations, translating term formulas.
In the displays below \(\Sigma\) denotes this relational presentation.
For every relational formula
\(\varphi(v_0,\ldots,v_{n-1})\), introduce a relation symbol whose
interpretation is its truth set, retaining equality. The resulting
definitional relational expansion is \(\Sigma^*\). In this presentation
\(T=\Sigma^*\), and \(\mathcal N(T)\) is the class of induced
subsystems of models of \(T\). Subsystems need not themselves be models.
The empty parameter system is permitted for type spaces.

A \emph{monomorphism} preserves and reflects every relation and is
injective. In the expanded language it is precisely a partial elementary
map of the original structures. Thus unions and amalgams below preserve
the original elementary information, not merely atomic information in
the unexpanded language. As in the source, \(|A|\) denotes the
universe of the system \(A\), whereas \(\kappa(X)\) denotes the
cardinality of a set \(X\), and \(\kappa(A)=\kappa(|A|)\) is the
cardinality, or power, of the system. A bare \(\kappa\) denotes a
cardinal; the argument distinguishes the cardinality operation
\(\kappa(\cdot)\). The source's equality of cardinals modulo
\(\aleph_0\) means equality after adding \(\aleph_0\).
The ambient set theory includes Choice.
\end{hbconvention}

\begin{hbtheorem}{thm:p03-morley-normalization}{proved=annex:p3.morley.normalization}{Definitional expansion, chains, and amalgamation}
Every \(\Sigma\)-model has exactly one expansion to a \(T\)-model.
Isomorphisms, partial elementary maps, model cardinalities, saturation,
and categoricity transfer in both directions. Every expanded formula is
equivalent modulo \(T\) to an atomic relation formula.
The theory \(T\) is complete. The class \(\mathcal N(T)\) is closed
under increasing unions, has joint embedding, and has amalgamation over
every common subsystem. Increasing unions of \(T\)-models are models.
\end{hbtheorem}

If an initial theory is categorical in an infinite cardinal and has no
finite models, completeness follows: two completions would, by
L\"owenheim--Skolem, give non-elementarily-equivalent models in that
cardinal. If finite models are allowed, the conclusion applies to its
infinite-model part, obtained by adding the sentences asserting at least
\(n\) elements for each finite \(n\). Finite models must not be folded
into the source's standing infinite complete theory.

\begin{hbdefinition}{def:p03-morley-type-space}{Formula algebras, points, and restriction maps}
For \(A\in\mathcal N(T)\), let \(T(A)\) be \(T\) together with the
diagram naming every element of \(A\). Because the language is
Morleyized, this is a complete parameter theory. Let \(F(A)\) be the
Boolean algebra of formulas with at most the free variable \(v_0\),
modulo equivalence in \(T(A)\). Its Stone space \(S(A)\) consists of
complete one-element types over \(A\), with clopen basis
\[
 U_\varphi=\{p\in S(A):\varphi\in p\}.
\]
A point is realized by \(b\) in an extension when every formula in it
holds of \(b\). For a monomorphism \(f:A\to B\), substitute
\(f(a)\) for named parameters to obtain the Boolean embedding
\(F(A)\to F(B)\). Its dual is the restriction map
\(f^*:S(B)\to S(A)\). For an inclusion we write \(i_{AB}^*\).
These maps, with all the spaces \(S(A)\), constitute the source's
category of continuous onto maps. Inverse images always refer to
actual restriction maps, not arbitrary continuous maps.
\end{hbdefinition}

\begin{hblemma}{thm:p03-morley-types}{proved=annex:p3.morley.types}{Realization and topology of parameter types}
Every \(S(A)\) is compact Hausdorff and zero-dimensional, with a
clopen basis of size at most \(\kappa(A)+\aleph_0\).
Every point is realized in an extension, and one extension realizes
all points at once. All restriction maps are continuous onto maps.
Two elements in possibly different extensions have the same type over
\(A\) exactly when the map identifying them and fixing \(A\) is a
monomorphism on the induced systems. Each basic clopen is the inverse
image of a basic clopen over a finite parameter subsystem.
\end{hblemma}

\subsection{The global derivative, rank, and finite degree}
\label{sec:p03-morley-global-derivative}

\begin{hbdefinition}{def:p03-morley-global-rank}{Morley's derivative across all extensions}
Simultaneously for all \(A\in\mathcal N(T)\), put
\[
 S^\alpha(A)=S(A)\setminus\bigcup_{\beta<\alpha}
                  \operatorname{Tr}^\beta(A).
\]
A point \(p\in S^\alpha(A)\) belongs to
\(\operatorname{Tr}^\alpha(A)\) when, for every monomorphism
\(f:A\to B\), every point of
\[
       (f^*)^{-1}(p)\cap S^\alpha(B)
\]
is isolated in \(S^\alpha(B)\). The rank of such a point is
\(\alpha\); points of \(\operatorname{Tr}^0(A)\) are called
\emph{algebraic}. A point surviving every stage has rank \(\infty\).
An isolated point of \(S(A)\) need not be algebraic: isolation in the
single space is weaker than isolation of every extension fiber.
The identity map is included among the tested maps.
\end{hbdefinition}

\begin{hbtheorem}{thm:p03-morley-derivative}{proved=annex:p3.morley.derivative}{Closedness, exact image, and finite localization}
For every ordinal \(\alpha\), \(S^\alpha(A)\) is closed and compact.
For every monomorphism \(f:A\to B\),
\[
 f^*(S^\alpha(B))=S^\alpha(A).
\]
For \(p\in S^\alpha(A)\),
\[
 p\in\operatorname{Tr}^\alpha(A)
 \quad\Longleftrightarrow\quad
 (f^*)^{-1}(p)\cap S^\alpha(B)
           \subseteq\operatorname{Tr}^\alpha(B).
\]
If \(p\) has ordinal rank \(\alpha\), some finite \(E\subseteq A\)
has \(i_{EA}^*(p)\) of exactly that rank.
\end{hbtheorem}

\begin{hbtheorem}{thm:p03-morley-degree}{proved=annex:p3.morley.degree}{Uniform finite multiplicity and its additive law}
If \(p\) has rank \(\alpha\), there is a positive integer bounding
\(\kappa((f^*)^{-1}(p)\cap S^\alpha(B))\) for every
\(f:A\to B\). The maximum of these integers is its degree
\(d(p)\). For every such map,
\[
 d(p)=\sum_{q\in(f^*)^{-1}(p)\cap S^\alpha(B)}d(q).
\]
In particular, an extension of the same rank has no larger degree;
if it has the same degree, it is the only extension of that rank.
Degree is a finite multiplicity, not the ordinal rank or the
cardinality of a basis.
\end{hbtheorem}

\begin{hbdefinition}{def:p03-morley-total-transcendence}{Stabilization and total transcendence}
Let \(\alpha_T\) be the least ordinal at which every global derivative
has stabilized: \(S^{\alpha_T}(A)=S^\beta(A)\) for every
\(\beta\geq\alpha_T\) and every \(A\). The theory is
\emph{totally transcendental} when this common residual space is empty.
In that case every type over every subsystem has ordinal rank.
\end{hbdefinition}

\begin{hbtheorem}{thm:p03-morley-stabilization}{proved=annex:p3.morley.stabilization}{Uniform stabilization and the countable rank bound}
The ordinal \(\alpha_T\) exists and is smaller than
\((2^{\aleph_0})^+\). If \(S^\alpha(A)\) is empty for one
subsystem, it is empty for every subsystem. The first empty stage is a
successor ordinal. If \(T\) is totally transcendental, then
\(\alpha_T<\omega_1\).
\end{hbtheorem}

\begin{hbtheorem}{thm:p03-morley-countability}{proved=annex:p3.morley.countability}{The original countability characterization}
If \(T\) is totally transcendental,
\[
 \kappa(S(A))\leq\kappa(A)+\aleph_0,
 \qquad
 \kappa(S(A))+\aleph_0=\kappa(A)+\aleph_0.
\]
Conversely, \(T\) is totally transcendental if \(S(A)\) is countable
for every countable \(A\in\mathcal N(T)\). If it is not totally
transcendental, some countable subsystem has \(2^{\aleph_0}\) types.
For the original countable theory this is equivalently
\(\omega\)-stability: countably many complete one-types over every
countable parameter set in a model.
\end{hbtheorem}

\begin{hbdefinition}{def:p03-morley-definable-rank}{Definable-set notation and its scope}
For a definable set \(X\) in finitely many variables, with arbitrary
parameters allowed, define \(\operatorname{RM}(X)\geq0\) when
\(X\ne\varnothing\); define
\(\operatorname{RM}(X)\geq\alpha+1\) when some elementary extension
has infinitely many pairwise disjoint definable subsets of \(X\),
each of rank at least \(\alpha\); and at a limit require every
earlier inequality. The empty set has rank \(-\infty\); a nonempty
set has its exact ordinal rank when the next inequality fails,
and rank \(\infty\) if none fails. Equivalently one works in a
sufficiently saturated ambient extension and allows its parameters.
For ordinal rank \(\alpha\), its Morley degree is the largest
number of disjoint definable subsets of rank \(\alpha\).
\end{hbdefinition}

\begin{hbtheorem}{thm:p03-morley-rank-bridge}{proved=annex:p3.morley.rank-bridge}{The bridge to definable-set Morley rank}
In a totally transcendental theory, the rank of a one-variable
\(\varphi\)-set over \(A\) is the largest rank of a point of
\(U_\varphi\subseteq S(A)\). Its degree is the sum of the degrees
of its points of that largest rank. These are exactly the splitting
rank and finite degree just defined. The same construction applies
to every finite tuple space. For countable theories, all definable
sets have ordinal Morley rank exactly when the theory is
\(\omega\)-stable. A rank-zero definable set is nonempty finite;
its degree is its number of elements. Rank and degree are not
algebraic-closure dimension.
\end{hbtheorem}

\subsection{The original three examples}
\label{sec:p03-morley-original-examples}

\begin{hbexample}{def:p03-morley-field-data}{Algebraically closed fields of a fixed characteristic}
Let \(\Sigma=\mathrm{ACF}_p\), where \(p=0\) or a fixed prime, and
let \(K(A)\) be the field generated by a parameter subsystem \(A\).
Over \(A\), the algebraic one-types correspond to irreducible
polynomials over \(K(A)\); there is one further type, asserting
nonvanishing of every nonzero polynomial over \(K(A)\).
This last type is the transcendental type. Morley's original example
uses characteristic zero; the same calculation works for each fixed
positive characteristic.
\end{hbexample}

\begin{hbtheorem}{thm:p03-morley-field-example}{proved=annex:p3.morley.field-example}{Field ranks, degrees, and the cardinal distinction}
The polynomial type for an irreducible polynomial has rank zero
and degree equal to its number of distinct roots (its separable
degree); in characteristic zero this is its polynomial degree.
The transcendental type has rank
one and degree one. Thus \(S^1(A)\) consists of that single type,
\(S^2(A)=\varnothing\), and \(\alpha_T=2\).
Every uncountable algebraically closed field of fixed characteristic
has transcendence degree equal to its cardinality; hence these
fields are categorical in every uncountable cardinal. Their
countably infinite models have the distinct transcendence degrees
\(0,1,2,\ldots,\aleph_0\), so they are not
\(\aleph_0\)-categorical. This numerical dimension is not the
one-variable rank one.
\end{hbtheorem}

\begin{hbexample}{def:p03-morley-tuple-code-data}{A categorical theory with larger finite rank}
Fix \(n\geq1\). There is an infinite unary part \(R_0\), and every
element outside it codes exactly one ordered \(n\)-tuple of distinct
elements of \(R_0\). Conversely each such tuple has exactly one
code, recorded by a relation \(R_1\) of arity \(n+1\).
For a parameter system \(B\), close it under taking code coordinates
and under forming codes; call the resulting set \(\overline B\).
\end{hbexample}

\begin{hbtheorem}{thm:p03-morley-tuple-code-example}{proved=annex:p3.morley.tuple-code-example}{Ranks in the tuple-code example}
This theory is categorical in every infinite cardinal. Elements of
\(\overline B\) have algebraic types. There is one type for a fresh
\(R_0\)-element, of rank one. A code type with exactly \(m\) of
its coordinates in \(\overline B\) has rank \(n-m\), for
\(m<n\). Consequently \(\alpha_T=n+1\). The calculation uses
extensions naming additional coordinates; counting isolated points
only in the initial space would miss these ranks.
\end{hbtheorem}

\begin{hbexample}{def:p03-morley-unary-data}{Closed subsets of Cantor space}
For a nonempty closed \(Y\subseteq2^\omega\), use unary predicates
\(R_j\) for \(j<\omega\). A finite Boolean predicate cell is empty
if its corresponding cylinder misses \(Y\), and infinite if that
cylinder meets \(Y\). These axioms give a complete theory: finite
formulas distinguish only finitely many predicate cells and equality
patterns. Write \(Y^{(\beta)}\) for the ordinary Cantor--Bendixson
derivatives, and \(\alpha_Y\) for their first stabilization stage.
\end{hbexample}

\begin{hbtheorem}{thm:p03-morley-unary-example}{proved=annex:p3.morley.unary-example}{The ordinary derivative inside the global derivative}
The named elements of \(A\) give the algebraic points of \(S(A)\).
The remaining points, one fresh-element type for each predicate
pattern in \(Y\), form \(S^1(A)\cong Y\). Restriction maps induce
homeomorphisms between these residual spaces, and
\[
 S^{1+\beta}(A)\cong Y^{(\beta)},\qquad
 \alpha_T=1+\alpha_Y.
\]
Here \(1+\beta\) is ordinal addition. The theory is totally
transcendental exactly when \(Y\) is countable, equivalently when
its perfect kernel is empty.
\end{hbtheorem}

\subsection{Ramsey models and total transcendence from categoricity}
\label{sec:p03-morley-ramsey-models}

\begin{hbdefinition}{def:p03-morley-skolem-hull}{Ordered indiscernibles and generated hulls}
An ordered set \((X,<)\) is a set of order indiscernibles in a
structure when increasing tuples of the same finite length satisfy
the same formulas. A countable Skolem expansion \(\Sigma^\#\)
is a conservative expansion by functions with universal axioms:
its reducts are \(\Sigma\)-models and every \(\Sigma\)-model
has an expansion. Its hull \(M(X,A)\) consists of all values of
terms on finite tuples from \(X\), including constant terms.
For elementary rather than merely relational preservation the
construction first Morleyizes and then Skolemizes, using
\(\Sigma^{*\#}\).
\end{hbdefinition}

\begin{hbtheorem}{thm:p03-morley-ramsey-em}{proved=annex:p3.morley.ramsey-em}{Ramsey, ordered indiscernibles, and the term-hull construction}
Every finite coloring of the \(n\)-element subsets of an infinite
set has an infinite homogeneous subset. Consequently, for any
theory with an infinite model and any ordered set \(X\), there is
a model containing \(X\) as order indiscernibles. In a Skolem
expansion, an order-preserving injection of the generators induces
a well-defined injective map of their term hull; a bijection induces
an automorphism. With the preliminary Morleyization these maps are
elementary on the original reduct. The hull has cardinality
\(\kappa(X)\) when \(X\) is infinite and the original
language is countable.
\end{hbtheorem}

\begin{hbtheorem}{thm:p03-morley-thin-models}{proved=annex:p3.morley.thin-models}{Thin models and the consequence of uncountable categoricity}
For every infinite cardinal \(\kappa\), a countable theory with
an infinite model has a model of size \(\kappa\) realizing only
countably many one-types over each countable subset of itself.
If the theory is categorical in some \(\kappa>\aleph_0\), it is
totally transcendental, hence \(\omega\)-stable.
Countable categoricity alone does not imply this: dense linear
order without endpoints has continuum many cut types over its
countable rational model.
\end{hbtheorem}

\begin{hbdefinition}{def:p03-morley-connected}{Connected and antisymmetric relations}
An \(n\)-ary relation is \emph{connected} on \(X\) if each
tuple of \(n\) distinct elements has a permutation satisfying it.
It is \emph{antisymmetric} in Morley's sense if each such tuple
also has a permutation not satisfying it. The second condition
is not the usual binary order axiom called antisymmetry.
\end{hbdefinition}

\begin{hbtheorem}{thm:p03-morley-no-order}{proved=annex:p3.morley.no-order}{No infinite permutation-sensitive relation}
In a totally transcendental theory no relation is both connected
and antisymmetric on an infinite subset of a model. The same holds
for a definable relation with a finite tuple of parameters.
\end{hbtheorem}

\subsection{Prime models compatible with chains}
\label{sec:p03-morley-prime-chains}

\begin{hbdefinition}{def:p03-morley-prime}{Prime extensions and unordered indiscernibles}
A model \(B\supseteq A\) is \emph{prime over} \(A\) when every
monomorphism of \(A\) into any model extends to a monomorphism of
\(B\) into that model. An \emph{indiscernible set over} \(A\)
is a set \(X\) disjoint from \(A\) for which every injective map
of \(A\cup X\) into itself fixing \(A\) is a monomorphism.
There is no distinguished order in this second notion.
\end{hbdefinition}

\begin{hbtheorem}{thm:p03-morley-prime-models}{proved=annex:p3.morley.prime-models}{Density and existence of prime extensions}
A subsystem \(A\) is a model exactly when its realized points
are dense in \(S(A)\). In a totally transcendental theory the
isolated points of each \(S(A)\) are dense: a minimum-rank point
of a nonempty open set is isolated in the whole space.
More generally, if these isolated points are dense for every
subsystem, every subsystem has a prime model over it, of size at
most \(\kappa(A)+\aleph_0\).
\end{hbtheorem}

\begin{hblemma}{thm:p03-morley-chain-types}{proved=annex:p3.morley.chain-types}{Coherent types and isolated lifts at limits}
Let \((A_\xi)_{\xi<\gamma}\) be an increasing nonempty chain,
with union \(A\), in a totally transcendental theory. A coherent
family \(p_\xi\in S(A_\xi)\) has eventually constant rank and
degree; it has a unique union type in \(S(A)\), of exactly those
eventual values. If \(p_0\) is isolated, it has coherent lifts
\(p_\xi\) that are isolated at every stage. One can include the
union stage; its lift is isolated there too.
\end{hblemma}

\begin{hbtheorem}{thm:p03-morley-chain-primes}{proved=annex:p3.morley.chain-primes}{Continuous chains of prime models}
For every increasing chain \((A_\xi)_{\xi<\gamma}\) in
\(\mathcal N(T)\) that is continuous at limit indices, there are
increasing models \((B_\xi)_{\xi<\gamma}\) with \(B_\xi\)
prime over \(A_\xi\), and
\(B_\delta=\bigcup_{\xi<\delta}B_\xi\) at limits.
The construction keeps each newly adjoined element isolated over
every later base together with the elements already adjoined.
This last condition is the transfer mechanism; separately chosen
prime models without it do not prove the theorem.
\end{hbtheorem}

\begin{hbtheorem}{thm:p03-morley-indiscernible-extraction}{proved=annex:p3.morley.indiscernible-extraction}{Large indiscernible sets inside the given system}
Suppose \(A\subseteq B\) are subsystems and
\(\kappa(A)<\kappa(B)=\kappa>\aleph_0\).
If \(\kappa\) is regular, \(B\setminus A\) contains an
indiscernible set over \(A\) of size \(\kappa\).
For singular \(\kappa\), for every \(\lambda<\kappa\) it
contains such a set of size greater than \(\lambda\).
This result assumes total transcendence and extracts its elements
inside \(B\), not merely inside an unspecified extension.
\end{hbtheorem}

\subsection{Saturation, omission transfer, and the final theorem}
\label{sec:p03-morley-saturation-transfer}

\begin{hbdefinition}{def:p03-morley-saturation}{Full and countable-subsystem saturation}
An infinite \(B\in\mathcal N(T)\) is saturated when every type
over each \(A\subseteq B\) with \(\kappa(A)<\kappa(B)\)
is realized in \(B\). For uncountable \(B\), saturation over
countable subsystems requires this only for countable \(A\).
Either property implies that \(B\) is a model. These are
realization properties, not an assertion that \(B\) realizes
every type over itself.
\end{hbdefinition}

\begin{hbtheorem}{thm:p03-morley-saturation-existence}{proved=annex:p3.morley.saturation-existence}{Uniqueness and the countable-saturation comparison model}
Two saturated models of the same infinite cardinality are
isomorphic. If \(T\) is totally transcendental and
\(\kappa>\aleph_0\), it has a model of cardinality \(\kappa\)
saturated over countable subsystems, with no assumption
\(\kappa^{\aleph_0}=\kappa\).
\end{hbtheorem}

\begin{hblemma}{thm:p03-morley-countable-omission-witness}{proved=annex:p3.morley.countable-omission-witness}{A countable witness with indiscernibles}
If an uncountable model \(B\) of a totally transcendental theory
is not saturated, there are a countable elementary submodel
\(A\subseteq B\), a subsystem \(A'\subseteq A\), an infinite
indiscernible set \(Y\subseteq A\setminus A'\) over \(A'\),
and a type \(q\in S(A')\) omitted in \(A\).
\end{hblemma}

\begin{hbtheorem}{thm:p03-morley-omission-transfer}{proved=annex:p3.morley.omission-transfer}{Transfer of a saturation failure to every uncountable size}
If a totally transcendental theory has any uncountable nonsaturated
model, then for each \(\kappa>\aleph_0\) it has a model of
cardinality \(\kappa\) that omits a type over a countable
subsystem. In particular that model is not saturated over
countable subsystems.
\end{hbtheorem}

\begin{hbtheorem}{thm:p03-morley-original-categoricity}{proved=annex:p3.morley.original-categoricity}{Morley's categoricity theorem and its stronger conclusion}
For a complete theory in a countable language, with an infinite
model, categoricity in one uncountable cardinal implies that
every uncountable model is saturated. Consequently it is
categorical in every uncountable cardinal. No continuum
hypothesis is required.
\end{hbtheorem}

The two comparison models in the categorical cardinal are an actual
countable-subsystem-saturated model and an actual model omitting a
countable-parameter type. Full saturation is concluded only after
the omission-transfer theorem rules out all uncountable failures.
This prevents replacing the hard transfer argument by the much
weaker observation that a suitably saturated extension exists.

\subsection{The original closing questions and the boundary of this source}
\label{sec:p03-morley-closing-questions}

The final page records questions rather than additional theorems
\cite[pp.~537--538]{Morley1965}. Historically, Morley asks whether an
uncountably categorical but not countably categorical theory has
exactly countably many countable models; whether such a theory can
be finitely axiomatizable; whether a language of size at most
\(\kappa>\aleph_0\) has categoricity transfer above \(\kappa\);
and whether a theory all of whose models have size at least
\(\kappa\) can be categorical in \(\kappa\).
For totally transcendental theories he asks for a model generated,
in the relevant symmetry sense, by a set all of whose permutations
extend to automorphisms; his footnote reports Silver's affirmative
answer. He also asks whether \(\alpha_T\) can ever be uncountable
(the proved total-transcendence case is negative), and what
conditions force it to be finite. He suggests categoricity and
finite axiomatizability as candidates, not as proved conditions.

These are the source's historical questions, not a declaration
that their 1965 status remains current. The later Baldwin--Lachlan
analysis and its countable-model consequences belong to their
own source presentation. The present proof uses neither an
unsupplied strongly minimal set nor an unsupplied prime-over-basis
theorem. The bibliographic works listed in Morley's references
retain their cited roles; that list does not allocate their whole
contents to this article's reconstruction.
